%% 第3巻 巻末（こたえ・記号カード・定数と単位・公式・チェック表・つぎの巻へ・さくいん・奥付）（第12回編集委員会）
%% 記号カードの一覧・記号を書くときのちゅうい・定数と単位・公式・表計算・つぎの巻へを、第3巻の本文に合わせて書きかえた（組 V3-9）
\makeatletter
%% ---------------- こたえ ----------------
\bkAnswerPages
\par\vspace{6mm}
\begin{tcolorbox}[bk base, colframe=bkTheme!60, colback=bkTheme!5, boxrule=0.8pt, arc=2mm, left=4mm, right=4mm, top=2.5mm, bottom=3mm]
{\sffamily\bfseries\color{bkDeep}\R{答}{こた}え\R{合}{あ}わせのしかた}\par\vspace{1.5mm}
\begin{enumerate}[label={\bkcircled{\arabic*}}, leftmargin=2.4\zw, itemsep=1pt, topsep=0pt]
\item まず\R{自分}{じ|ぶん}の\R{答}{こた}えを\R{全部}{ぜん|ぶ}\R{書}{か}いてから、ここを\R{見}{み}よう。
\item ちがっていたら、\R{答}{こた}えだけを\R{直}{なお}さずに、\R{同}{おな}じ\R{形}{かたち}の\R{例題}{れい|だい}にもどって、どの\R{行}{ぎょう}でちがったかをさがそう。
\item \R{単位}{たん|い}もいっしょに\R{確}{たし}かめよう。\R{数}{かず}が\R{合}{あ}っていても、\R{単位}{たん|い}がちがえば、ちがう\R{量}{りょう}です（2-1）。
\end{enumerate}
\end{tcolorbox}

%% ---------------- 記号カード ----------------
%% 前の巻の記号で、この巻でも使うもの（この巻の章の \bkkigocard は自動でうしろに集まる）
\bkKigoPages{%
\bkkigorow{$=$}{イコール}{\R{左}{ひだり}と\R{右}{みぎ}が\R{同}{おな}じ\R{大}{おお}きさ}{1-4}%
\bkkigorow{$-$}{マイナス}{ひき\R{算}{ざん}。0より\R{小}{ちい}さい\R{数}{かず}の\R{印}{しるし}にも\R{使}{つか}う（2-2）}{1-5}%
\bkkigorow{$\times$}{かける}{1つ\R{分}{ぶん}の\R{数}{かず}×いくつ\R{分}{ぶん}}{1-12}%
\bkkigorow{$\div$}{わる}{\R{同}{おな}じ\R{数}{かず}ずつ\R{分}{わ}ける。1あたりを\R{求}{もと}める}{1-14}%
\bkkigorow{$\dfrac{1}{2}$}{2\R{分}{ぶん}の1}{1を2つに\R{等}{ひと}しく\R{分}{わ}けた1つ\R{分}{ぶん}}{1-16}%
\bkkigorow{$<$\ \ $>$\ \ $\leqq$\ \ $\geqq$}{\R{不等号}{ふ|とう|ごう}}{\R{大}{おお}きさをくらべる\R{記号}{き|ごう}。$\leqq$ は「\R{以下}{い|か}」、$\geqq$ は「\R{以上}{い|じょう}」。\R{変域}{へん|いき}（3-1）や\R{不等式}{ふ|とう|しき}（3-5）で\R{使}{つか}う}{2-2}%
\bkkigorow{$\Delta$}{デルタ}{\R{変化量}{へん|か|りょう}。「あと－まえ」で\R{求}{もと}める。$\Delta x$ は $x$ の\R{変化量}{へん|か|りょう}}{2-3}%
\bkkigorow{\%}{パーセント}{\R{百分率}{ひゃく|ぶん|りつ}の\R{記号}{き|ごう}。\R{割合}{わり|あい}の0.01が1\%}{2-5}%
\bkkigorow{$\pm$}{プラスマイナス}{ずれのはば。3-10 の $x=\pm3$ は「3と $-3$ の2つ」の\R{印}{しるし}}{2-6}%
\bkkigorow{$2^3$}{2の3\R{乗}{じょう}}{2を3\R{回}{かい}かけた\R{数}{かず}。$2\times2\times2=8$}{2-7}%
\bkkigorow{$\sqrt{\phantom{00}}$}{ルート}{2\R{乗}{じょう}すると\R{中}{なか}の\R{数}{かず}になる\R{正}{せい}の\R{数}{かず}。$\sqrt{9}=3$}{2-8}%
\bkkigorow{$\pi$}{パイ}{\R{円周率}{えん|しゅう|りつ}。\R{円周}{えん|しゅう}÷\R{直径}{ちょっ|けい}。$3.14159\cdots$}{2-9}%
\bkkigorow{$x$}{エックス}{\R{数}{かず}のかわりに\R{使}{つか}う\R{文字}{も|じ}。この\R{巻}{かん}では\R{変数}{へん|すう}（3-1）}{2-11}%
\bkkigorow{$(\ \ )$}{かっこ}{ひとまとまりにする\R{印}{しるし}。かっこの\R{中}{なか}を\R{先}{さき}に\R{計算}{けい|さん}する}{2-11}%
\bkkigorow{$x^{-1}$}{エックスのマイナス1\R{乗}{じょう}}{$\dfrac{1}{x}$ のこと。$x$ で1\R{回}{かい}わる}{2-12}%
}

%% 記号カードのつづき：手で書くときのちゅうい（この本は問題をノートに解く）
\def\bkbkhead#1{\par\addvspace{4mm}\noindent{\sffamily\bfseries\large\color{\bkTT}#1}\par\nobreak\vspace{2mm}}
\clearpage
\bkbkhead{\R{記号}{き|ごう}を\R{書}{か}くときのちゅうい}
{\small この\R{巻}{かん}では、\R{文字}{も|じ}の\R{右下}{みぎ|した}や\R{右上}{みぎ|うえ}に\R{小}{ちい}さく\R{書}{か}く\R{記号}{き|ごう}や、ギリシャ\R{文字}{も|じ}がふえます。\R{手}{て}で\R{書}{か}くと、ほかの\R{字}{じ}や\R{記号}{き|ごう}とまぎれやすいものの、\R{書}{か}くときのコツです。}\par\vspace{2mm}
\begin{bkcenter}
\begin{tblr}{colspec={Q[c,m,wd=18mm] Q[l,m,wd=24mm] Q[l,m,wd=84mm]}, row{1}={bg=bkVol!14, font=\sffamily\bfseries\footnotesize}, hlines={0.4pt,bkLine}, rowsep=1.6pt, cells={font=\small}}
\R{記号}{き|ごう} & まぎれやすいもの & \R{書}{か}くときのコツ \\
%% 第14回 A5：記号と単位が並ぶときの書き方を1行足した
$V$〔V〕 & $V$\,V（VV） & \R{記号}{き|ごう}と\R{単位}{たん|い}が\R{並}{なら}ぶときは、\R{単位}{たん|い}を〔\ 〕でくくる。「\R{電圧}{でん|あつ} $V$〔V〕」「\R{電流}{でん|りゅう} $I$〔A〕」（3-1） \\
$f(x)$ & $f\times x$ & かっこは「$x$ を\R{入}{い}れる」\R{印}{しるし}で、かけ\R{算}{ざん}ではない。$f$ のすぐ\R{右}{みぎ}に\R{書}{か}く（3-3） \\
$x_1$ & $x1$、$x^1$ & 1は\R{右下}{みぎ|した}に\R{小}{ちい}さく。\R{右上}{みぎ|うえ}に\R{書}{か}くと\R{累乗}{るい|じょう}になってしまう（3-2） \\
$\lvert x\rvert$ & 1、$l$（エル） & たての\R{線}{せん}を、\R{数}{かず}の\R{両側}{りょう|がわ}に\R{同}{おな}じ\R{長}{なが}さで\R{書}{か}く（3-2） \\
$\Delta y$ & $\Delta\times y$ & $\Delta y$ で1つの\R{量}{りょう}（$y$ の\R{変化量}{へん|か|りょう}）。$\Delta$ と $y$ をはなさない（3-5） \\
$x^{\frac{1}{2}}$ & $x\dfrac{1}{2}$ & \R{分数}{ぶん|すう}の\R{指数}{し|すう}も、\R{右上}{みぎ|うえ}に\R{小}{ちい}さく\R{書}{か}く（3-11） \\
$\log_2 8$ & $\log 28$ & \R{底}{てい}の2は\R{右下}{みぎ|した}に\R{小}{ちい}さく。8はふつうの\R{大}{おお}きさで\R{書}{か}く（3-12） \\
$\alpha$ & $a$ & \R{丸}{まる}から\R{右}{みぎ}へ、2\R{本}{ほん}の\R{線}{せん}を\R{開}{ひら}いてのばす。$a$ のようにたての\R{線}{せん}を\R{書}{か}かない（3-10） \\
$\beta$ & $B$ & \R{左}{ひだり}のたての\R{線}{せん}を、\R{下}{した}へ\R{長}{なが}くのばす（3-10） \\
$\pm$ & $+$ & $+$ の\R{下}{した}に\R{横線}{よこ|せん}を1\R{本}{ぽん}（3-10） \\
$\bar{x}$ & $x$ & \R{横}{よこ}ぼうを $x$ のすぐ\R{上}{うえ}に。$\bar{x}$（\R{平均}{へい|きん}）と $x$ はちがう\R{数}{かず}（3-13） \\
$\sigma$ & 6、$b$、$\Sigma$ & \R{丸}{まる}の\R{上}{うえ}から\R{右}{みぎ}へ\R{横}{よこ}ぼうをのばす。\R{大文字}{おお|も|じ}の $\Sigma$ は\R{別}{べつ}の\R{記号}{き|ごう}（5-10）（3-14） \\
${}_n\mathrm{C}_r$ & $n\mathrm{C}r$ & $n$ は\R{左下}{ひだり|した}、$r$ は\R{右下}{みぎ|した}に\R{小}{ちい}さく。C は\R{大}{おお}きく\R{書}{か}く（3-17） \\
$P(A)$ & $p$ & \R{確率}{かく|りつ}の $P$ は\R{大文字}{おお|も|じ}。\R{圧力}{あつ|りょく}の $p$（3-7）と\R{区別}{く|べつ}する（3-17） \\
\end{tblr}
\end{bkcenter}

%% ---------------- 定数・単位・公式 ----------------
\bkbacktitle{よく\R{使}{つか}う\R{定数}{てい|すう}・\R{単位}{たん|い}・\R{公式}{こう|しき}}
\bksubsection{1}{この\R{巻}{かん}で\R{使}{つか}う\R{単位}{たん|い}}
\begin{bkcenter}
\small
\begin{tblr}{colspec={Q[l,m,wd=22mm] Q[l,m,wd=30mm] Q[l,m,wd=57mm] Q[c,m,wd=12mm]}, row{1}={bg=bkVol!14, font=\sffamily\bfseries\footnotesize}, hlines={0.4pt,bkLine}, rowsep=1.4pt}
\R{量}{りょう} & \R{単位}{たん|い}（\R{読}{よ}み\R{方}{かた}） & ほかの\R{単位}{たん|い}との\R{関係}{かん|けい} & \R{章}{しょう} \\
\R{時間}{じ|かん} & s（\R{秒}{びょう}）、min（\R{分}{ふん}）、h（\R{時間}{じ|かん}） & 1\,min＝60\,s、1\,h＝60\,min & {3-6\\3-11} \\
\R{速}{はや}さ & m/s、m/min、km/h & \R{道}{みち}のり÷\R{時間}{じ|かん}。$x$-$t$ グラフの\R{傾}{かたむ}き & {3-4\\3-5} \\
\R{傾}{かたむ}き・\R{変化}{へん|か}の\R{速}{はや}さ & ℃/h など & \R{傾}{かたむ}きの\R{単位}{たん|い}は、たての\R{単位}{たん|い}÷\R{横}{よこ}の\R{単位}{たん|い} & 3-5 \\
{\R{速}{はや}さの\\ふえ\R{方}{かた}} & m/s$^2$ & 1\R{秒}{びょう}ごとに、\R{速}{はや}さが\R{何}{なん}\,m/s ふえるか。$g$ の\R{単位}{たん|い} & 3-1 \\
\R{力}{ちから}・ばね\R{定数}{てい|すう} & N、N/m & ばね\R{定数}{てい|すう}は、1\,m のばすのに\R{何}{なん}\,N いるか（$F=kx$） & {3-1\\3-16} \\
\R{圧力}{あつ|りょく} & {Pa（パスカル）\\kPa} & 1\,m$^2$ あたり1\,N の\R{力}{ちから}が1\,Pa。$\qty{1}{kPa}=\qty{1000}{Pa}$ & 3-7 \\
\R{温度}{おん|ど} & ℃（\R{摂氏}{せっ|し}）、°F（\R{華氏}{か|し}） & $F=1.8C+32$ & 3-3 \\
\R{体積}{たい|せき} & L、mL & $\qty{1}{L}=\qty{1000}{mL}$ & 3-7 \\
{\R{電圧}{でん|あつ}・\R{電流}{でん|りゅう}\\\R{抵抗}{てい|こう}} & V、A、Ω & オームの\R{法則}{ほう|そく} $V=RI$。$\qty{1}{A}=\qty{1000}{mA}$ & {3-6\\3-12} \\
\R{電気}{でん|き}の\R{使用量}{し|よう|りょう} & kWh（キロワット\R{時}{じ}） & \R{電気料金}{でん|き|りょう|きん}の\R{計算}{けい|さん}に\R{使}{つか}う & 3-3 \\
\R{強}{つよ}さの\R{比}{ひ} & dB（デシベル） & \R{強}{つよ}さが10\R{倍}{ばい}で10\,dB、100\R{倍}{ばい}で20\,dB & 3-12 \\
\R{振動数}{しん|どう|すう} & Hz（ヘルツ） & 100\,Hz は、1\R{秒}{びょう}に100\R{回}{かい}のゆれ & 3-12 \\
\end{tblr}
\end{bkcenter}

\bksubsection{2}{よく\R{使}{つか}う\R{数}{かず}}
\begin{bkcenter}
\small
\begin{tblr}{colspec={Q[l,m,wd=48mm] Q[l,m,wd=56mm] Q[c,m,wd=16mm]}, row{1}={bg=bkVol!14, font=\sffamily\bfseries\footnotesize}, hlines={0.4pt,bkLine}, rowsep=1.4pt}
\R{名前}{な|まえ} & \R{値}{あたい} & \R{章}{しょう} \\
\R{落}{お}ちる\R{物}{もの}の\R{速}{はや}さのふえ\R{方}{かた} $g$ & \R{約}{やく}\qty{9.8}{m/s^2}（\R{地球}{ち|きゅう}の\R{表面}{ひょう|めん}の\R{近}{ちか}く） & 3-1 \\
\R{真空}{しん|くう}の\R{中}{なか}の\R{光}{ひかり}の\R{速}{はや}さ $c$ & \R{約}{やく} $3.0\times10^8$\,m/s（ちょうど\qty{299792458}{m/s}） & 3-1 \\
\R{円周率}{えん|しゅう|りつ} $\pi$ & 3.14159…（\R{数学}{すう|がく}の\R{定数}{てい|すう}） & {2-9\\3-1} \\
$\sqrt{2}$ & 1.41421…（$=2^{\frac{1}{2}}$） & {2-8\\3-11} \\
$\sqrt{3}$ & 1.73205… & 3-10 \\
$\dfrac{1}{\sqrt{2}}=\dfrac{\sqrt{2}}{2}$ & 0.70710… & 3-10 \\
$\log_{10}2$、\ $\log_{10}3$ & 0.301…、\ 0.477… & 3-12 \\
\R{大気圧}{たい|き|あつ} & \R{約}{やく}\qty{101}{kPa}（1013\,hPa） & 3-7 \\
\end{tblr}
\end{bkcenter}

\bksubsection{3}{\R{関数}{かん|すう}とグラフ・\R{連立方程式}{れん|りつ|ほう|てい|しき}}
{\small
\begin{longtblr}{colspec={Q[l,m,wd=29mm] Q[l,m,wd=80mm] Q[c,m,wd=15mm]}, rowhead=1, row{1}={bg=bkVol!14, font=\sffamily\bfseries\footnotesize}, hlines={0.4pt,bkLine}, rowsep=1.6pt}
\R{公式}{こう|しき}・きまり & \R{式}{しき} & \R{章}{しょう} \\
\R{物理}{ぶつ|り}の\R{式}{しき} & \R{落}{お}ちる\R{物}{もの}の\R{速}{はや}さ $v=gt$、\ \R{重力}{じゅう|りょく} $W=mg$、\ ばね $F=kx$ & 3-1 \\
\R{点}{てん}を\R{動}{うご}かす & \R{点}{てん} $(a,\ b)$ を\R{右}{みぎ}へ $p$、\R{上}{うえ}へ $q$ \R{動}{うご}かすと $(a+p,\ b+q)$ & 3-2 \\
2\R{点}{てん}の\R{距離}{きょ|り} & \R{横}{よこ}に\R{並}{なら}んだ2\R{点}{てん}は $\lvert x_2-x_1\rvert$、たてに\R{並}{なら}んだ2\R{点}{てん}は $\lvert y_2-y_1\rvert$ & 3-2 \\
\R{関数}{かん|すう} $f(x)$ & $f(x)=x^2$ なら $f(3)=9$、\ $f(a+h)=(a+h)^2$（かっこごと\R{入}{い}れる） & 3-3 \\
\R{合成関数}{ごう|せい|かん|すう} & $g(x)=3x+1$、$f(u)=u^2$ なら $f(g(x))=(3x+1)^2$（\R{内側}{うち|がわ}から\R{計算}{けい|さん}） & 3-3 \\
\R{摂氏}{せっ|し}と\R{華氏}{か|し} & $F=1.8C+32$ & 3-3 \\
\R{傾}{かたむ}き & $\dfrac{\Delta y}{\Delta x}=\dfrac{y_2-y_1}{x_2-x_1}$（たての\R{変化量}{へん|か|りょう}÷\R{横}{よこ}の\R{変化量}{へん|か|りょう}） & 3-5 \\
1\R{次関数}{じ|かん|すう} & $y=ax+b$（$a$ が\R{傾}{かたむ}き＝\R{変化}{へん|か}の\R{割合}{わり|あい}、$b$ が\R{切片}{せっ|ぺん}） & 3-5 \\
1\R{点}{てん}と\R{傾}{かたむ}きから & \R{傾}{かたむ}き $m$、\R{点}{てん} $(x_1,\ y_1)$ を\R{通}{とお}る\R{直線}{ちょく|せん}：\ $y-y_1=m(x-x_1)$ & 3-5 \\
グラフと\R{速}{はや}さ & $x$-$t$ グラフの\R{傾}{かたむ}き＝\R{速}{はや}さ。$v$-$t$ グラフの\R{面積}{めん|せき}＝\R{道}{みち}のり & 3-5 \\
1\R{次不等式}{じ|ふ|とう|しき} & \R{両辺}{りょう|へん}にマイナスの\R{数}{かず}をかける・わると、\R{不等号}{ふ|とう|ごう}の\R{向}{む}きが\R{反対}{はん|たい}になる & 3-5 \\
\R{連立方程式}{れん|りつ|ほう|てい|しき} & \R{代入法}{だい|にゅう|ほう}：「$y=\cdots$」をもう1つの\R{式}{しき}にかっこごと\R{入}{い}れる。\R{加減法}{か|げん|ほう}：\R{係数}{けい|すう}をそろえて、\R{式}{しき}どうしをたす・ひく & 3-6 \\
\R{解}{かい}の\R{個数}{こ|すう} & 2\R{本}{ほん}の\R{直線}{ちょく|せん}が\R{交}{まじ}わる→1\R{組}{くみ}。\R{平行}{へい|こう}→\R{解}{かい}なし。\R{重}{かさ}なる→\R{無数}{む|すう} & 3-6 \\
\R{反比例}{はん|ぴ|れい} & $y=\dfrac{a}{x}$（$xy=a$）。グラフは\R{双曲線}{そう|きょく|せん} & 3-7 \\
ボイルの\R{法則}{ほう|そく} & \R{温度}{おん|ど}が\R{変}{か}わらなければ、\R{気体}{き|たい}の $pV$ は\R{一定}{いっ|てい} & 3-7 \\
\end{longtblr}\par}

\bksubsection{4}{\R{式}{しき}の\R{計算}{けい|さん}・2\R{次関数}{じ|かん|すう}・2\R{次方程式}{じ|ほう|てい|しき}}
{\small
\begin{longtblr}{colspec={Q[l,m,wd=29mm] Q[l,m,wd=80mm] Q[c,m,wd=15mm]}, rowhead=1, row{1}={bg=bkVol!14, font=\sffamily\bfseries\footnotesize}, hlines={0.4pt,bkLine}, rowsep=1.6pt}
\R{公式}{こう|しき}・きまり & \R{式}{しき} & \R{章}{しょう} \\
\R{展開}{てん|かい} & $(a+b)(c+d)=ac+ad+bc+bd$（\R{全部}{ぜん|ぶ}の\R{組}{く}み\R{合}{あ}わせをかけてたす） & 3-8 \\
2\R{乗}{じょう}の\R{展開}{てん|かい} & $(a+b)^2=a^2+2ab+b^2$、\ $(a-b)^2=a^2-2ab+b^2$ & 3-8 \\
\R{和}{わ}と\R{差}{さ}の\R{積}{せき} & $(a+b)(a-b)=a^2-b^2$ & 3-8 \\
$(x+a)(x+b)$ & $(x+a)(x+b)=x^2+(a+b)x+ab$ & 3-8 \\
3\R{乗}{じょう}の\R{展開}{てん|かい} & $(a+b)^3=a^3+3a^2b+3ab^2+b^3$ & 3-8 \\
\R{恒等式}{こう|とう|しき} & \R{文字}{も|じ}にどんな\R{数}{かず}を\R{入}{い}れても\R{成}{な}り\R{立}{た}つ\R{等式}{とう|しき}。\R{係数}{けい|すう}を\R{比}{くら}べる・\R{数}{かず}を\R{代入}{だい|にゅう}する、で\R{係数}{けい|すう}を\R{決}{き}める & 3-8 \\
\R{平方完成}{へい|ほう|かん|せい} & $x^2+bx+c=(x+p)^2-p^2+c$（$p$ は $b$ の\R{半分}{はん|ぶん}）。\ \R{例}{れい}：$x^2+6x+10=(x+3)^2+1$ & {3-8\\3-9} \\
\R{放物線}{ほう|ぶつ|せん} & $y=ax^2$：$a>0$ なら\R{上}{うえ}に\R{開}{ひら}く、$a<0$ なら\R{下}{した}に\R{開}{ひら}く & 3-9 \\
\R{放物線}{ほう|ぶつ|せん}を\R{動}{うご}かす & $y=a(x-p)^2+q$：\R{頂点}{ちょう|てん} $(p,\ q)$、\R{軸}{じく}は\R{直線}{ちょく|せん} $x=p$ & 3-9 \\
$x^2=k$ & $k>0$ なら $x=\pm\sqrt{k}$。$k=0$ なら $x=0$。$k<0$ なら\R{解}{かい}はない & 3-10 \\
かけて0 & $A\times B=0$ なら、$A=0$ または $B=0$ & 3-10 \\
\R{因数分解}{いん|すう|ぶん|かい} & $x^2-a^2=(x+a)(x-a)$。$x^2+\fbox{\R{和}{わ}}\,x+\fbox{\R{積}{せき}}$ は、たして\R{和}{わ}・かけて\R{積}{せき}の2つの\R{数}{かず}をさがす & 3-10 \\
$\sqrt{\phantom{2}}$ の\R{計算}{けい|さん} & $\sqrt{12}=\sqrt{4}\times\sqrt{3}=2\sqrt{3}$、\ $\dfrac{1}{\sqrt{2}}=\dfrac{\sqrt{2}}{2}$（\R{有理化}{ゆう|り|か}） & 3-10 \\
\R{解}{かい}の\R{公式}{こう|しき} & $ax^2+bx+c=0$ の\R{解}{かい}は\quad $x=\dfrac{-b\pm\sqrt{b^2-4ac}}{2a}$ & 3-10 \\
\R{判別式}{はん|べつ|しき} & $D=b^2-4ac$。$D>0$ なら\R{解}{かい}は2つ、$D=0$ なら1つ、$D<0$ ならない & 3-10 \\
\R{解}{かい}と\R{因数分解}{いん|すう|ぶん|かい} & \R{解}{かい}が $\alpha$、$\beta$ なら\quad $ax^2+bx+c=a(x-\alpha)(x-\beta)$ & 3-10 \\
\end{longtblr}\par}

\bksubsection{5}{\R{指数}{し|すう}と\R{対数}{たい|すう}}
{\small
\begin{longtblr}{colspec={Q[l,m,wd=29mm] Q[l,m,wd=80mm] Q[c,m,wd=15mm]}, rowhead=1, row{1}={bg=bkVol!14, font=\sffamily\bfseries\footnotesize}, hlines={0.4pt,bkLine}, rowsep=1.6pt}
\R{公式}{こう|しき}・きまり & \R{式}{しき} & \R{章}{しょう} \\
\R{倍々}{ばい|ばい}でふえる・へる & （$x$ \R{回後}{かい|ご}の\R{量}{りょう}）＝（はじめの\R{量}{りょう}）×$a^x$（$a$ は1\R{回}{かい}ごとの\R{倍率}{ばい|りつ}） & 3-11 \\
\R{指数}{し|すう}のきまり & $a^{x+h}=a^x\times a^h$、\ $(a^m)^n=a^{mn}$、\ $a^0=1$、\ $a^{-x}=\dfrac{1}{a^x}$ & 3-11 \\
\R{分数}{ぶん|すう}の\R{指数}{し|すう} & $x^{\frac{1}{2}}=\sqrt{x}$、\ $x^{-\frac{1}{2}}=\dfrac{1}{\sqrt{x}}$ & 3-11 \\
\R{半減期}{はん|げん|き} $T$ & （のこっている\R{量}{りょう}）＝（はじめの\R{量}{りょう}）×$\left(\dfrac{1}{2}\right)^{\frac{t}{T}}$ & 3-11 \\
\R{対数}{たい|すう} & $a^p=b$ と $\log_ab=p$ は\R{同}{おな}じこと。\ $a^{\log_ab}=b$ & 3-12 \\
\R{対数}{たい|すう}のきまり & $\log_aMN=\log_aM+\log_aN$、\ $\log_a\dfrac{M}{N}=\log_aM-\log_aN$、\ $\log_aM^k=k\log_aM$（$a>0$、$a\neq1$、$M>0$、$N>0$。\R{底}{てい} $a$ が\R{何}{なん}でも\R{成}{な}り\R{立}{た}つ） & 3-12 \\
\R{特別}{とく|べつ}な\R{値}{あたい} & $\log_a1=0$、\ $\log_aa=1$ & 3-12 \\
\R{底}{てい}を\R{変}{か}える & $\log_ab=\dfrac{\log_{10}b}{\log_{10}a}$（\R{電卓}{でん|たく}の log は\R{底}{てい}が10） & 3-12 \\
\R{対数}{たい|すう}のグラフ & $y=\log_ax$ は、$y=a^x$ を\R{直線}{ちょく|せん} $y=x$ で\R{折}{お}り\R{返}{かえ}したもの & 3-12 \\
デシベル & （dB）＝10×$\log_{10}$（\R{電力}{でん|りょく}の\R{比}{ひ}）＝20×$\log_{10}$（\R{電圧}{でん|あつ}の\R{比}{ひ}） & 3-12 \\
\end{longtblr}\par}

\bksubsection{6}{データと\R{確率}{かく|りつ}}
{\small
\begin{longtblr}{colspec={Q[l,m,wd=29mm] Q[l,m,wd=80mm] Q[c,m,wd=15mm]}, rowhead=1, row{1}={bg=bkVol!14, font=\sffamily\bfseries\footnotesize}, hlines={0.4pt,bkLine}, rowsep=1.6pt}
\R{公式}{こう|しき}・きまり & \R{式}{しき} & \R{章}{しょう} \\
\R{平均}{へい|きん} & $\bar{x}=\dfrac{x_1+x_2+\cdots+x_n}{n}$ & 3-13 \\
\R{中央値}{ちゅう|おう|ち} & \R{小}{ちい}さい\R{順}{じゅん}に\R{並}{なら}べて、まん\R{中}{なか}の\R{値}{あたい}（\R{偶数}{ぐう|すう}こなら、まん\R{中}{なか}の2つの\R{平均}{へい|きん}） & 3-13 \\
{\R{相対度数}{そう|たい|ど|すう}\\\R{範囲}{はん|い}} & {\R{相対度数}{そう|たい|ど|すう}＝\R{度数}{ど|すう}÷\R{合計}{ごう|けい}\\\R{範囲}{はん|い}＝\R{最大値}{さい|だい|ち}－\R{最小値}{さい|しょう|ち}} & 3-13 \\
\R{四分位範囲}{し|ぶん|い|はん|い} & \R{第}{だい}3\R{四分位数}{し|ぶん|い|すう}－\R{第}{だい}1\R{四分位数}{し|ぶん|い|すう} & 3-13 \\
\R{偏差}{へん|さ} & \R{値}{あたい}－\R{平均}{へい|きん}。\R{全部}{ぜん|ぶ}たすと、かならず0 & 3-14 \\
\R{分散}{ぶん|さん} & \R{偏差}{へん|さ}の2\R{乗}{じょう}の\R{平均}{へい|きん}（\R{合計}{ごう|けい}を\R{個数}{こ|すう} $n$ でわる） & 3-14 \\
\R{標準偏差}{ひょう|じゅん|へん|さ} & $\sigma=\sqrt{\dfrac{(x_1-\bar{x})^2+(x_2-\bar{x})^2+\cdots+(x_n-\bar{x})^2}{n}}$ & 3-14 \\
$\sigma$ の\R{性質}{せい|しつ} & \R{全部}{ぜん|ぶ}に $b$ をたしても\R{同}{おな}じ。$a$ \R{倍}{ばい}（$a>0$）すると $a$ \R{倍}{ばい} & 3-14 \\
\R{平均}{へい|きん}のずれ & $n$ \R{回}{かい}はかった\R{平均}{へい|きん}のずれは $\dfrac{\sigma}{\sqrt{n}}$ & 3-14 \\
\R{共分散}{きょう|ぶん|さん} & {$\sigma_{xy}=\{(x_1-\bar{x})(y_1-\bar{y})+(x_2-\bar{x})(y_2-\bar{y})$\\\hspace*{3em}$+\cdots+(x_n-\bar{x})(y_n-\bar{y})\}\div n$} & 3-15 \\
\R{相関係数}{そう|かん|けい|すう} & $r=\dfrac{\sigma_{xy}}{\sigma_x\,\sigma_y}$。いつも $-1\leqq r\leqq1$ & 3-15 \\
{\R{最小二乗法}{さい|しょう|に|じょう|ほう}\\（\R{原点}{げん|てん}を\R{通}{とお}る）} & $y=ax$ の\quad $a=$（$xy$ の\R{和}{わ}）÷（$x^2$ の\R{和}{わ}） & 3-16 \\
\R{最小二乗法}{さい|しょう|に|じょう|ほう} & $y=ax+b$ の\quad $a=\dfrac{\sigma_{xy}}{\sigma_x^{\,2}}$（\R{共分散}{きょう|ぶん|さん}÷$x$ の\R{分散}{ぶん|さん}）、\ $b=\bar{y}-a\bar{x}$ & 3-16 \\
\R{確率}{かく|りつ} & $P(A)=\dfrac{a}{n}$\quad（$A$ の\R{起}{お}こる\R{場合}{ば|あい}の\R{数}{かず}）÷（\R{全部}{ぜん|ぶ}の\R{場合}{ば|あい}の\R{数}{かず}） & 3-17 \\
\R{組}{くみ}\R{合}{あ}せの\R{数}{かず} & ${}_5\mathrm{C}_2=\dfrac{5\times4}{2\times1}=10$。パスカルの\R{三角形}{さん|かく|けい}の\R{数}{かず}と\R{同}{おな}じ & 3-17 \\
\R{期待値}{き|たい|ち} & （\R{値}{あたい}×その\R{確率}{かく|りつ}）を\R{全部}{ぜん|ぶ}たした\R{数}{かず}。サイコロの\R{目}{め}なら3.5 & 3-17 \\
\R{正規分布}{せい|き|ぶん|ぷ} & \R{平均}{へい|きん}±$\sigma$ に\R{約}{やく}68\%、\R{平均}{へい|きん}±$2\sigma$ に\R{約}{やく}95\% が\R{入}{はい}る & {3-14\\3-17} \\
\end{longtblr}\par}

\bksubsection{7}{\R{表計算}{ひょう|けい|さん}・\R{電卓}{でん|たく}・プログラムの\R{書}{か}き\R{方}{かた}}
{\small
\begin{longtblr}{colspec={Q[l,m,wd=25mm] Q[l,m,wd=47mm] Q[l,m,wd=42mm] Q[c,m,wd=12mm]}, rowhead=1, row{1}={bg=bkVol!14, font=\sffamily\bfseries\footnotesize}, hlines={0.4pt,bkLine}, rowsep=1.6pt, colsep=4pt}
\R{求}{もと}めるもの & \R{表計算}{ひょう|けい|さん}・\R{電卓}{でん|たく}での\R{書}{か}き\R{方}{かた}（\R{例}{れい}） & \R{意味}{い|み} & \R{章}{しょう} \\
0でわった\R{印}{しるし} & \texttt{\#DIV/0!} & 0でわったというエラー & 3-7 \\
\R{累乗}{るい|じょう} & {\texttt{=2\^{}0.5}\\\R{電卓}{でん|たく}は \texttt{x\^{}y} のキー} & $2^{0.5}$（$=\sqrt{2}$） & 3-11 \\
\R{常用対数}{じょう|よう|たい|すう} & {\texttt{=LOG10(2)}\\\R{電卓}{でん|たく}は \texttt{log} のキー} & $\log_{10}2$ & 3-12 \\
\R{平均}{へい|きん} & \texttt{=AVERAGE(B2:B13)} & $\bar{x}$ & 3-13 \\
\R{中央値}{ちゅう|おう|ち} & \texttt{=MEDIAN(B2:B13)} & まん\R{中}{なか}の\R{値}{あたい} & 3-13 \\
\R{四分位数}{し|ぶん|い|すう} & \texttt{=QUARTILE.INC(B2:B13,3)} & \R{第}{だい}3\R{四分位数}{し|ぶん|い|すう} & 3-13 \\
\R{平方根}{へい|ほう|こん} & \texttt{=SQRT(D23)} & $\sqrt{\phantom{2}}$ & 3-14 \\
\R{標準偏差}{ひょう|じゅん|へん|さ} & \texttt{=STDEV.P(B2:B21)} & $\sigma$（$n$ でわる） & 3-14 \\
\R{共分散}{きょう|ぶん|さん} & {\texttt{=COVARIANCE.P(}\\\texttt{\ \ A2:A11,B2:B11)}} & $\sigma_{xy}$ & 3-15 \\
\R{相関係数}{そう|かん|けい|すう} & \texttt{=CORREL(A2:A11,B2:B11)} & $r$ & 3-15 \\
\R{傾}{かたむ}き & \texttt{=SLOPE(B2:B8,A2:A8)} & \R{最小二乗法}{さい|しょう|に|じょう|ほう}の $a$ & 3-16 \\
\R{切片}{せっ|ぺん} & \texttt{=INTERCEPT(B2:B8,A2:A8)} & \R{最小二乗法}{さい|しょう|に|じょう|ほう}の $b$ & 3-16 \\
\R{乱数}{らん|すう} & \texttt{=RANDBETWEEN(0,1)} & 0か1を、でたらめに\R{出}{だ}す & 3-17 \\
{プログラムの\\\R{関数}{かん|すう}（Python）} & {\texttt{def f(x):}\\\texttt{\ \ \ \ return 1.8 * x + 32}} & \R{箱}{はこ} \texttt{f} のきまりを\R{決}{き}める。\texttt{return} は\R{計算}{けい|さん}した\R{数}{かず}を\R{出}{だ}す & 3-3 \\
\end{longtblr}\par}
\begin{itemize}[leftmargin=1.4\zw, itemsep=2pt, topsep=4pt, label={\color{\bkTT}●}]
\item \texttt{STDEV.P} は\R{偏差}{へん|さ}の2\R{乗}{じょう}の\R{合計}{ごう|けい}を $n$ でわり、その\R{平方根}{へい|ほう|こん}（$\sqrt{\phantom{0}}$）をとります。\texttt{STDEV.S} は $n$ のかわりに $n-1$ でわります。この\R{本}{ほん}では $n$ でわります（3-14）。
\item \texttt{SLOPE} と \texttt{INTERCEPT} のかっこの\R{中}{なか}は、はじめに $y$ の\R{範囲}{はん|い}、\R{次}{つぎ}に $x$ の\R{範囲}{はん|い}の\R{順}{じゅん}です（3-16）。
\item \R{電卓}{でん|たく}の \texttt{log} のキーは、\R{底}{てい}が10の\R{常用対数}{じょう|よう|たい|すう}です。$\log_2$ は、\R{底}{てい}を\R{変}{か}える\R{式}{しき}で\R{常用対数}{じょう|よう|たい|すう}になおして\R{求}{もと}めます（3-12）。
\end{itemize}

%% ---------------- チェック表 ----------------
\bkbacktitle{できるようになったことチェック\R{表}{ひょう}}
{\small \R{各章}{かく|しょう}の「ふりかえり」を\R{集}{あつ}めた\R{表}{ひょう}です。\R{章}{しょう}を\R{学}{まな}び\R{終}{お}えたら1\R{回}{かい}め、\R{次}{つぎ}の\R{巻}{かん}に\R{進}{すす}む\R{前}{まえ}に2\R{回}{かい}めのチェックをしよう。できなかった\R{所}{ところ}は、ページにもどって\R{読}{よ}みなおします。}\par\vspace{2mm}
%% 自動生成（tools/mkchecklist.py）。各章の「ふりかえり」の項目を集めたもの
\begingroup\footnotesize
\def\bkcb{\tikz[baseline=-0.5ex]\draw[\bkTF, line width=0.7pt, rounded corners=0.5pt] (0,-0.3em) rectangle (0.8em,0.5em);}
\begin{longtblr}{colspec={Q[c,m,wd=9mm] X[l,m] Q[c,m,wd=10mm] Q[c,m,wd=9mm] Q[c,m,wd=9mm]}, width=\linewidth, rowhead=1, colsep=3pt,
  row{1}={bg=bkVol!14, font=\sffamily\bfseries\scriptsize}, hlines={0.4pt,bkLine}, rowsep=2.3pt}
\R{章}{しょう} & できるようになったこと & ページ & 1\R{回}{かい}め & 2\R{回}{かい}め \\
{\sffamily\bfseries\color{\bkTT}3-1} & \R{式}{しき}の\R{中}{なか}の\R{変数}{へん|すう}と\R{定数}{てい|すう}を\R{見分}{み|わ}けられる。 & \pageref{bk:3-1:hensu} & \bkcb & \bkcb \\
 & \R{物理定数}{ぶつ|り|てい|すう}と、\R{問題}{もん|だい}ごとに\R{決}{き}まる\R{数}{かず}のちがいを\R{説明}{せつ|めい}できる。 & \pageref{bk:3-1:butsuri} & \bkcb & \bkcb \\
 & \R{変域}{へん|いき}を\R{不等号}{ふ|とう|ごう}で\R{書}{か}ける。 & \pageref{bk:3-1:heniki} & \bkcb & \bkcb \\
{\sffamily\bfseries\color{\bkTT}3-2} & \R{原点}{げん|てん}と\R{軸}{じく}を\R{決}{き}めて、\R{点}{てん}をかいたり\R{座標}{ざ|ひょう}を\R{読}{よ}んだりできる。 & \pageref{bk:3-2:zahyo} & \bkcb & \bkcb \\
 & マイナスの\R{座標}{ざ|ひょう}、4つの\R{区画}{く|かく}、\R{点}{てん}を\R{動}{うご}かした\R{座標}{ざ|ひょう}がわかる。 & \pageref{bk:3-2:ugokasu} & \bkcb & \bkcb \\
 & \R{絶対値}{ぜっ|たい|ち}を\R{使}{つか}って、2\R{点}{てん}の\R{距離}{きょ|り}を\R{求}{もと}められる。 & \pageref{bk:3-2:kyori} & \bkcb & \bkcb \\
{\sffamily\bfseries\color{\bkTT}3-3} & \R{関数}{かん|すう}を\R{見分}{み|わ}け、\R{表}{ひょう}・グラフ・\R{式}{しき}で\R{表}{あらわ}せる。 & \pageref{bk:3-3:kansu} & \bkcb & \bkcb \\
 & $f(x)$ を\R{読}{よ}み、$f(3)$ や $f(a+h)$ を\R{計算}{けい|さん}できる。 & \pageref{bk:3-3:fx} & \bkcb & \bkcb \\
 & 2つの\R{関数}{かん|すう}をつないだ\R{合成}{ごう|せい}\R{関数}{かん|すう}を\R{作}{つく}れる。 & \pageref{bk:3-3:gousei} & \bkcb & \bkcb \\
{\sffamily\bfseries\color{\bkTT}3-4} & \R{軸}{じく}と\R{目}{め}もりを\R{決}{き}めて、はかった\R{点}{てん}からグラフをかける。 & \pageref{bk:3-4:jiku} & \bkcb & \bkcb \\
 & $x$-$t$ グラフと $v$-$t$ グラフを\R{読}{よ}み\R{分}{わ}けられる。 & \pageref{bk:3-4:xt} & \bkcb & \bkcb \\
 & \R{表}{ひょう}にない\R{値}{あたい}を\R{線}{せん}から\R{読}{よ}み、\R{外挿}{がい|そう}の\R{注意}{ちゅう|い}を\R{説明}{せつ|めい}できる。 & \pageref{bk:3-4:naisou} & \bkcb & \bkcb \\
{\sffamily\bfseries\color{\bkTT}3-5} & \R{傾}{かたむ}きを「たての\R{変化量}{へん|か|りょう}÷\R{横}{よこ}の\R{変化量}{へん|か|りょう}」で\R{求}{もと}め、\R{正}{せい}・\R{負}{ふ}・0がわかる。 & \pageref{bk:3-5:katamuki} & \bkcb & \bkcb \\
 & $y=ax+b$ の\R{傾}{かたむ}きと\R{切片}{せっ|ぺん}を\R{読}{よ}み、$x$-$t$ グラフの\R{傾}{かたむ}きを\R{速}{はや}さとして\R{使}{つか}える。 & \pageref{bk:3-5:ichiji} & \bkcb & \bkcb \\
 & 1\R{点}{てん}と\R{傾}{かたむ}きから、$y-y_1=m(x-x_1)$ で\R{直線}{ちょく|せん}の\R{式}{しき}を\R{作}{つく}れる。 & \pageref{bk:3-5:tenkatamuki} & \bkcb & \bkcb \\
{\sffamily\bfseries\color{\bkTT}3-6} & \R{交点}{こう|てん}が、2つの\R{式}{しき}を\R{同時}{どう|じ}にみたす\R{組}{くみ}だとわかる。 & \pageref{bk:3-6:kouten} & \bkcb & \bkcb \\
 & \R{代入法}{だい|にゅう|ほう}と\R{加減法}{か|げん|ほう}で\R{解}{と}き、\R{両方}{りょう|ほう}の\R{式}{しき}で\R{確}{たし}かめられる。 & \pageref{bk:3-6:dainyuu} & \bkcb & \bkcb \\
 & \R{解}{かい}がない・\R{無数}{む|すう}のわけを\R{説明}{せつ|めい}し、\R{出会}{で|あ}い・\R{買}{か}い\R{物}{もの}・\R{回路}{かい|ろ}の\R{問題}{もん|だい}を\R{解}{と}ける。 & \pageref{bk:3-6:heikou} & \bkcb & \bkcb \\
{\sffamily\bfseries\color{\bkTT}3-7} & \R{双曲線}{そう|きょく|せん}をかき、\R{軸}{じく}にくっつかないわけを\R{言}{い}える。 & \pageref{bk:3-7:jiku} & \bkcb & \bkcb \\
 & $a<0$ のグラフを\R{読}{よ}み、\R{変化}{へん|か}の\R{割合}{わり|あい}を\R{計算}{けい|さん}できる。 & \pageref{bk:3-7:wariai} & \bkcb & \bkcb \\
 & $xy$ が\R{一定}{いっ|てい}を\R{見}{み}つけ、\R{式}{しき}を\R{作}{つく}って\R{使}{つか}える。 & \pageref{bk:3-7:pv} & \bkcb & \bkcb \\
{\sffamily\bfseries\color{\bkTT}3-8} & かっこのある\R{式}{しき}を\R{展開}{てん|かい}し、\R{面積}{めん|せき}の\R{図}{ず}で\R{説明}{せつ|めい}できる。 & \pageref{bk:3-8:kumiawase} & \bkcb & \bkcb \\
 & \R{恒等式}{こう|とう|しき}の\R{係数}{けい|すう}を、2\R{通}{とお}りの\R{方法}{ほう|ほう}で\R{決}{き}められる。 & \pageref{bk:3-8:koutou} & \bkcb & \bkcb \\
 & \R{平方完成}{へい|ほう|かん|せい}して、いちばん\R{小}{ちい}さい・\R{大}{おお}きい\R{値}{あたい}を\R{読}{よ}める。 & \pageref{bk:3-8:heihou} & \bkcb & \bkcb \\
{\sffamily\bfseries\color{\bkTT}3-9} & 2\R{乗}{じょう}に\R{比例}{ひ|れい}を\R{見}{み}つけ、$y=ax^2$ の\R{式}{しき}とグラフにできる。 & \pageref{bk:3-9:nijo} & \bkcb & \bkcb \\
 & $a$ とグラフの\R{形}{かたち}の\R{関係}{かん|けい}や、\R{変化}{へん|か}の\R{割合}{わり|あい}を\R{説明}{せつ|めい}できる。 & \pageref{bk:3-9:a} & \bkcb & \bkcb \\
 & \R{平方完成}{へい|ほう|かん|せい}で\R{頂点}{ちょう|てん}を、\R{範囲}{はん|い}から\R{最大}{さい|だい}・\R{最小}{さい|しょう}を\R{求}{もと}められる。 & \pageref{bk:3-9:heihou} & \bkcb & \bkcb \\
{\sffamily\bfseries\color{\bkTT}3-10} & $x^2=k$ の\R{形}{かたち}や\R{因数分解}{いん|すう|ぶん|かい}で、2\R{次方程式}{じ|ほう|てい|しき}を\R{解}{と}ける。 & \pageref{bk:3-10:insu} & \bkcb & \bkcb \\
 & \R{解}{かい}の\R{公式}{こう|しき}を\R{作}{つく}り、$\sqrt{\phantom{2}}$ の\R{計算}{けい|さん}をして\R{解}{と}ける。 & \pageref{bk:3-10:kousiki} & \bkcb & \bkcb \\
 & \R{判別式}{はん|べつ|しき}で\R{解}{かい}の\R{個数}{こ|すう}を\R{調}{しら}べ、\R{場面}{ば|めん}に\R{合}{あ}う\R{解}{かい}を\R{選}{えら}べる。 & \pageref{bk:3-10:hanbetsu} & \bkcb & \bkcb \\
{\sffamily\bfseries\color{\bkTT}3-11} & 「ふえる\R{量}{りょう}が、いまある\R{量}{りょう}に\R{比例}{ひ|れい}する」\R{変化}{へん|か}を、\R{表}{ひょう}・\R{式}{しき}・グラフにできる。 & \pageref{bk:3-11:hirei} & \bkcb & \bkcb \\
 & \R{指数}{し|すう}のきまり $a^{x+h}=a^x\times a^h$、$(a^m)^n=a^{mn}$、$a^0=1$ を\R{使}{つか}える。 & \pageref{bk:3-11:kimari} & \bkcb & \bkcb \\
 & $x^{\frac{1}{2}}=\sqrt{x}$ のわけを\R{説明}{せつ|めい}できる。 & \pageref{bk:3-11:bunsu} & \bkcb & \bkcb \\
{\sffamily\bfseries\color{\bkTT}3-12} & 10を\R{何乗}{なん|じょう}するかでけたを\R{数}{かぞ}え、$\log_ab$ の\R{意味}{い|み}を\R{説明}{せつ|めい}できる。 & \pageref{bk:3-12:log} & \bkcb & \bkcb \\
 & \R{対数}{たい|すう}のきまりを、\R{指数}{し|すう}のきまりから\R{作}{つく}って\R{使}{つか}える。 & \pageref{bk:3-12:kimari} & \bkcb & \bkcb \\
 & \R{片対数}{かた|たい|すう}グラフやデシベルを\R{使}{つか}える。 & \pageref{bk:3-12:kata} & \bkcb & \bkcb \\
{\sffamily\bfseries\color{\bkTT}3-13} & \R{何回}{なん|かい}もはかった\R{値}{あたい}の\R{平均}{へい|きん}を、$\bar{x}$ を\R{使}{つか}って\R{求}{もと}められる。 & \pageref{bk:3-13:heikin} & \bkcb & \bkcb \\
 & \R{度数分布表}{ど|すう|ぶん|ぷ|ひょう}とヒストグラムを\R{作}{つく}り、\R{読}{よ}みとれる。 & \pageref{bk:3-13:dosuu} & \bkcb & \bkcb \\
 & \R{中央値}{ちゅう|おう|ち}・\R{最頻値}{さい|ひん|ち}・\R{四分位数}{し|ぶん|い|すう}を\R{求}{もと}め、\R{箱}{はこ}ひげ\R{図}{ず}をかける。 & \pageref{bk:3-13:daihyou} & \bkcb & \bkcb \\
{\sffamily\bfseries\color{\bkTT}3-14} & \R{偏差}{へん|さ}を\R{求}{もと}め、\R{合計}{ごう|けい}が0になるわけを\R{説明}{せつ|めい}できる。 & \pageref{bk:3-14:hensa} & \bkcb & \bkcb \\
 & \R{分散}{ぶん|さん}と $\sigma$ を\R{求}{もと}め、ばらつきをくらべられる。 & \pageref{bk:3-14:sigma} & \bkcb & \bkcb \\
 & \R{平均}{へい|きん}$\pm\sigma$ と $\tfrac{\sigma}{\sqrt{n}}$ の\R{意味}{い|み}を\R{説明}{せつ|めい}できる。 & \pageref{bk:3-14:nanawari} & \bkcb & \bkcb \\
{\sffamily\bfseries\color{\bkTT}3-15} & \R{散布図}{さん|ぷ|ず}をかき、\R{相関}{そう|かん}の\R{向}{む}きと\R{強}{つよ}さを\R{見分}{み|わ}けられる。 & \pageref{bk:3-15:sanpu} & \bkcb & \bkcb \\
 & \R{共分散}{きょう|ぶん|さん}を\R{計算}{けい|さん}し、その\R{符号}{ふ|ごう}の\R{意味}{い|み}を\R{説明}{せつ|めい}できる。 & \pageref{bk:3-15:kyobunsan} & \bkcb & \bkcb \\
 & \R{相関係数}{そう|かん|けい|すう} $r$ を\R{求}{もと}め、\R{相関}{そう|かん}と\R{因果}{いん|が}を\R{区別}{く|べつ}できる。 & \pageref{bk:3-15:r} & \bkcb & \bkcb \\
{\sffamily\bfseries\color{\bkTT}3-16} & ずれの2\R{乗}{じょう}の\R{和}{わ}で、どの\R{直線}{ちょく|せん}がよく\R{合}{あ}うかくらべられる。 & \pageref{bk:3-16:zure} & \bkcb & \bkcb \\
 & $y=ax$ の $a$ を、\R{平方完成}{へい|ほう|かん|せい}で\R{求}{もと}められる。 & \pageref{bk:3-16:genten} & \bkcb & \bkcb \\
 & $y=ax+b$ を\R{共分散}{きょう|ぶん|さん}と\R{分散}{ぶん|さん}で\R{求}{もと}め、\R{残差}{ざん|さ}で\R{確}{たし}かめられる。 & \pageref{bk:3-16:ippan} & \bkcb & \bkcb \\
{\sffamily\bfseries\color{\bkTT}3-17} & \R{表}{ひょう}や\R{樹形図}{じゅ|けい|ず}で\R{数}{かぞ}えて、\R{確率}{かく|りつ}と\R{期待値}{き|たい|ち}を\R{求}{もと}められる。 & \pageref{bk:3-17:kazoeru} & \bkcb & \bkcb \\
 & パスカルの\R{三角形}{さん|かく|けい}で、\R{表}{おもて}が $k$ \R{枚}{まい}になる\R{数}{かず}を\R{求}{もと}められる。 & \pageref{bk:3-17:pascal} & \bkcb & \bkcb \\
 & \R{正規分布}{せい|き|ぶん|ぷ}の68\%と95\%を\R{使}{つか}える。 & \pageref{bk:3-17:seiki} & \bkcb & \bkcb \\
\end{longtblr}
\endgroup


%% ---------------- つぎの巻へ（第4巻の予告と、この巻の道具がこの先どこで使われるか） ----------------
\bkbacktitle{つぎの\R{巻}{かん}へ}
\newcommand\bknextpill[1]{\bk@pill{#1}}
\def\bknextrow{1}
\newcommand\bknextcard[4]{%
  \begin{tcolorbox}[enhanced, nobeforeafter, box align=top, width=0.485\linewidth, colframe=bkVol4!75!black, colback=white, boxrule=0.9pt, arc=2.5mm,
    left=2.5mm, right=2.5mm, top=1.5mm, bottom=2mm, equal height group=bknext\bknextrow]
  \bknextpill{#1}\enspace{\sffamily\bfseries\small #2}\par\vspace{1mm}
  {\sffamily\bfseries\footnotesize\color{bkDeep4}#3}\par\vspace{1mm}
  {\footnotesize\setlength{\parindent}{0pt}{\sffamily\bfseries\color{bkInk}\R{使}{つか}う\R{道具}{どう|ぐ}}\enspace #4}
  \end{tcolorbox}}
\R{表}{ひょう}・グラフ・\R{式}{しき}で\R{変化}{へん|か}のきまりをつかむ\R{道具}{どう|ぐ}がそろいました。\R{第}{だい}4\R{巻}{かん}「\R{形}{かたち}と\R{向}{む}き」では、この\R{巻}{かん}で\R{身}{み}につけた\R{座標}{ざ|ひょう}・\R{関数}{かん|すう}・グラフを\R{使}{つか}って、\R{長}{なが}さと\R{角度}{かく|ど}で、\R{形}{かたち}と\R{向}{む}きを\R{数}{かず}にします。\R{第}{だい}4\R{巻}{かん}の\R{章}{しょう}をのぞいてみましょう。\par\vspace{2.5mm}
\noindent\bknextcard{4-2}{\R{角度}{かく|ど}と\R{円}{えん}、ラジアン}{\R{角度}{かく|ど}を「\R{長}{なが}さ」ではかる\R{方法}{ほう|ほう}}{\R{円周率}{えん|しゅう|りつ} $\pi$（3-1）と、1\R{次関数}{じ|かん|すう}・\R{比例}{ひ|れい}（3-5）。}\hfill
\bknextcard{4-4}{\R{三平方}{さん|へい|ほう}の\R{定理}{てい|り}}{ななめの\R{距離}{きょ|り}を、たてと\R{横}{よこ}から\R{計算}{けい|さん}する}{\R{座標}{ざ|ひょう}と2\R{点}{てん}の\R{距離}{きょ|り}（3-2）、\R{展開}{てん|かい}（3-8）、$x^2=k$ と $\sqrt{\phantom{2}}$ の\R{計算}{けい|さん}（3-10）。}\par\vspace{1.2mm}
\def\bknextrow{2}\noindent\bknextcard{4-6}{\R{坂道}{さか|みち}と\R{三角比}{さん|かく|ひ}}{\R{坂}{さか}の\R{上}{うえ}の\R{荷物}{に|もつ}にかかる\R{力}{ちから}を、2つに\R{分}{わ}ける}{\R{坂}{さか}の\R{急}{きゅう}さを\R{表}{あらわ}す\R{傾}{かたむ}き（3-5）と、\R{比}{ひ}（2-5）。}\hfill
\bknextcard{4-8}{\R{三角関数}{さん|かく|かん|すう}：\R{円運動}{えん|うん|どう}の\R{影}{かげ}}{\R{回}{まわ}る\R{点}{てん}の\R{影}{かげ}は、\R{行}{い}ったり\R{来}{き}たりする}{\R{座標}{ざ|ひょう}（3-2）、\R{関数}{かん|すう}（3-3）と、グラフの\R{読}{よ}み\R{方}{かた}（3-4・3-9）。}\par\vspace{1.2mm}
\def\bknextrow{3}\noindent\bknextcard{4-10}{\R{向}{む}きと\R{大}{おお}きさをもつ\R{量}{りょう}：ベクトル}{\R{2人}{ふたり}で\R{引}{ひ}くと、\R{荷物}{に|もつ}はどっちへ\R{動}{うご}く？}{\R{座標}{ざ|ひょう}と、\R{点}{てん}を\R{動}{うご}かす\R{計算}{けい|さん}（3-2）。}\hfill
\bknextcard{4-13}{\R{空間}{くう|かん}の\R{座標}{ざ|ひょう}と3\R{次元}{じ|げん}のベクトル}{\R{部屋}{へ|や}の\R{中}{なか}の1\R{点}{てん}を、3つの\R{数}{かず}で\R{表}{あらわ}す}{2つの\R{数}{かず}で\R{場所}{ば|しょ}を\R{表}{あらわ}す\R{座標}{ざ|ひょう}と、\R{下}{した}つき\R{文字}{も|じ}（3-2）。}\par
\bkfukidashi{delta}{\R{第}{だい}4\R{巻}{かん}からは、\R{漢字}{かん|じ}にふりがながつかないよ。チェック\R{表}{ひょう}で\R{自信}{じ|しん}のない\R{章}{しょう}を\R{見}{み}なおしてから、\R{第}{だい}4\R{巻}{かん}へ\R{進}{すす}もう！}

\clearpage
\bkbkhead{この\R{巻}{かん}の\R{道具}{どう|ぐ}は、この\R{先}{さき}でこう\R{使}{つか}う}
{\small \R{第}{だい}3\R{巻}{かん}の\R{章}{しょう}と、その\R{道具}{どう|ぐ}を\R{使}{つか}う\R{第}{だい}4\R{巻}{かん}から\R{先}{さき}の\R{章}{しょう}です（\R{章}{しょう}のはじめの「\R{道}{みち}しるべ」の「この\R{章}{しょう}の\R{先}{さき}に」から\R{集}{あつ}めました）。}\par\vspace{3mm}
\begingroup\let\bkRorig\R\def\R#1#2{\mbox{\bkRorig{#1}{#2}}}% 第7回：表の中でルビの語を途中で切らない
%% 自動生成（tools/mknext.py）。各章の道しるべ「この章の先に」から、後ろの巻の章を集めたもの
\begin{longtblr}{colspec={Q[l,m,wd=58mm] X[l,m]}, width=\linewidth, rowhead=1, row{1}={bg=bkVol!14, font=\sffamily\bfseries\footnotesize},
  hlines={0.4pt,bkLine}, rowsep=1.5pt, cells={font=\footnotesize}}
この\R{巻}{かん}の\R{章}{しょう} & この\R{道具}{どう|ぐ}を\R{使}{つか}う\R{章}{しょう} \\
{\bknextpill{3-1}\ \R{変}{か}わる\R{量}{りょう}と\R{変}{か}わらない\R{量}{りょう}} & {\R{第}{だい}4\R{巻}{かん}から\R{先}{さき}：\bknextpill{5-4}\,\bknextpill{5-5}} \\
{\bknextpill{3-2}\ \R{場所}{ば|しょ}を\R{数}{かず}で\R{表}{あらわ}す：\R{座標}{ざ|ひょう}} & {\R{第}{だい}4\R{巻}{かん}から\R{先}{さき}：\bknextpill{4-4}\,\bknextpill{4-10}\,\bknextpill{4-13}\,\bknextpill{4-15}} \\
{\bknextpill{3-3}\ \R{関数}{かん|すう}は「\R{入}{い}れると\R{出}{で}てくる」ルール} & {\R{第}{だい}4\R{巻}{かん}から\R{先}{さき}：\bknextpill{5-1}\,\bknextpill{5-3}\,\bknextpill{5-6}\,\bknextpill{9-8}} \\
{\bknextpill{3-4}\ グラフをかこう・\R{読}{よ}もう} & {\R{第}{だい}4\R{巻}{かん}から\R{先}{さき}：\bknextpill{5-1}\,\bknextpill{5-9}} \\
{\bknextpill{3-5}\ グラフの\R{傾}{かたむ}きは「\R{変化}{へん|か}の\R{速}{はや}さ」} & {\R{第}{だい}4\R{巻}{かん}から\R{先}{さき}：\bknextpill{5-1}\,\bknextpill{5-3}\,\bknextpill{8-2}} \\
{\bknextpill{3-6}\ 2\R{本}{ほん}の\R{直線}{ちょく|せん}が\R{交}{まじ}わるところ：\R{連立方程式}{れん|りつ|ほう|てい|しき}} & {\R{第}{だい}4\R{巻}{かん}から\R{先}{さき}：\bknextpill{7-1}\,\bknextpill{9-10}} \\
{\bknextpill{3-7}\ \R{反比例}{はん|ぴ|れい}のグラフ} & {\R{第}{だい}4\R{巻}{かん}から\R{先}{さき}：\bknextpill{5-4}\,\bknextpill{5-14}} \\
{\bknextpill{3-8}\ かっこをはずす：\R{展開}{てん|かい}と\R{平方完成}{へい|ほう|かん|せい}} & {\R{第}{だい}4\R{巻}{かん}から\R{先}{さき}：\bknextpill{5-4}\,\bknextpill{9-7}} \\
{\bknextpill{3-9}\ \R{放物線}{ほう|ぶつ|せん}と $y=ax^2$} & {\R{第}{だい}4\R{巻}{かん}から\R{先}{さき}：\bknextpill{5-1}\,\bknextpill{5-2}} \\
{\bknextpill{3-10}\ ボールはいつその\R{高}{たか}さに\R{届}{とど}く？：2\R{次方程式}{じ|ほう|てい|しき}} & {\R{第}{だい}4\R{巻}{かん}から\R{先}{さき}：\bknextpill{5-5}\,\bknextpill{5-7}\,\bknextpill{5-11}\,\bknextpill{6-12}} \\
{\bknextpill{3-11}\ \R{指数関数}{し|すう|かん|すう}：いまある\R{量}{りょう}に\R{比例}{ひ|れい}してふえる・へる} & {\R{第}{だい}4\R{巻}{かん}から\R{先}{さき}：\bknextpill{5-14}\,\bknextpill{5-18}\,\bknextpill{6-4}} \\
{\bknextpill{3-12}\ \R{対数}{たい|すう}：けたを\R{数}{かぞ}える\R{道具}{どう|ぐ}} & {\R{第}{だい}4\R{巻}{かん}から\R{先}{さき}：\bknextpill{5-14}\,\bknextpill{9-8}} \\
{\bknextpill{3-14}\ ばらつきを1つの\R{数}{かず}で：\R{標準偏差}{ひょう|じゅん|へん|さ}} & {\R{第}{だい}4\R{巻}{かん}から\R{先}{さき}：\bknextpill{9-4}} \\
{\bknextpill{3-15}\ 2つの\R{量}{りょう}の\R{関係}{かん|けい}：\R{散布図}{さん|ぷ|ず}と\R{相関}{そう|かん}} & {\R{第}{だい}4\R{巻}{かん}から\R{先}{さき}：\bknextpill{4-12}} \\
{\bknextpill{3-16}\ データに\R{直線}{ちょく|せん}を\R{引}{ひ}く：\R{最小二乗法}{さい|しょう|に|じょう|ほう}} & {\R{第}{だい}4\R{巻}{かん}から\R{先}{さき}：\bknextpill{5-7}\,\bknextpill{7-5}} \\
{\bknextpill{3-17}\ \R{偶然}{ぐう|ぜん}のばらつき：\R{確率}{かく|りつ}と\R{正規分布}{せい|き|ぶん|ぷ}} & {\R{第}{だい}4\R{巻}{かん}から\R{先}{さき}：\bknextpill{5-17}\,\bknextpill{9-4}} \\
\end{longtblr}

\endgroup

%% ---------------- さくいん ----------------
\bkbacktitle{さくいん}
{\small \R{数字}{すう|じ}はページです。その\R{語}{ご}がはじめて\R{出}{で}てくるページや、くわしく\R{説明}{せつ|めい}しているページをのせています。}\par\vspace{2mm}
\begin{multicols}{2}\small\setlength{\parskip}{0pt}
\directlua{bk.printidx()}
\end{multicols}

%% ---------------- 奥付 ----------------
\bkpadfour
\global\bk@decofalse\bkfolio{false}
\null\vfill
\noindent\begin{minipage}{\linewidth}\small
{\sffamily\bfseries\Large\color{bkDeep}\R{物理}{ぶつ|り}を\R{理解}{り|かい}するための\R{数学}{すう|がく}　\R{第}{だい}3\R{巻}{かん}}\par\vspace{1mm}
{\sffamily\bfseries\large \R{関数}{かん|すう}とデータ}\par\vspace{3mm}
{\color{bkVol}\hrule height 1pt}\vspace{3mm}
\begin{tabular}{@{}l@{\hspace{6mm}}l@{}}
\R{編者}{へん|しゃ} & 「\R{物理}{ぶつ|り}を\R{理解}{り|かい}するための\R{数学}{すう|がく}」\R{編集委員会}{へん|しゅう|い|いん|かい} \\
 & \R{委員長}{い|いん|ちょう}　\R{小学校長}{しょう|がっ|こう|ちょう}（まとめ\R{役}{やく}） \\
 & \R{委員}{い|いん}　\R{数学}{すう|がく}・\R{電気工学}{でん|き|こう|がく}・\R{機械工学}{き|かい|こう|がく}・\R{教育学}{きょう|いく|がく}・\R{量子力学}{りょう|し|りき|がく}・\R{建築工学}{けん|ちく|こう|がく}・\\
 & \phantom{\R{委員}{い|いん}　}\R{半導体工学}{はん|どう|たい|こう|がく}の\R{大学教授}{だい|がく|きょう|じゅ}、\R{幼稚園教諭}{よう|ち|えん|きょう|ゆ}、\R{塾講師}{じゅく|こう|し}、\\
 & \phantom{\R{委員}{い|いん}　}\R{特別支援学校教諭}{とく|べつ|し|えん|がっ|こう|きょう|ゆ}、\R{心理士}{しん|り|し}、\R{科学解説者}{か|がく|かい|せつ|しゃ}、\R{小学校教諭}{しょう|がっ|こう|きょう|ゆ} \\
デザイン & ユニバーサルデザイン\R{担当委員}{たん|とう|い|いん} \\
\R{挿絵}{さし|え} & イラスト\R{担当委員}{たん|とう|い|いん} \\
\R{版}{はん} & \R{第}{だい}1\R{版}{はん}（\R{第}{だい}12\R{回編集委員会}{かい|へん|しゅう|い|いん|かい}）　2026\R{年}{ねん}10\R{月}{がつ} \\
\end{tabular}\par\vspace{3mm}
{\color{bkVol}\hrule height 1pt}\vspace{2mm}
{\footnotesize この\R{本}{ほん}の\R{問題}{もん|だい}の\R{答}{こた}えは、\R{巻末}{かん|まつ}の「こたえ」にあります。くわしい\R{解}{と}き\R{方}{かた}は、\R{別冊}{べっ|さつ}の\R{問題集}{もん|だい|しゅう}と\R{指導書}{し|どう|しょ}にのせます。}
\end{minipage}
\makeatother
